\documentclass[11pt]{article}
\usepackage[a4paper,margin=27mm]{geometry}
\usepackage{amsmath,amssymb,amsthm,xcolor,booktabs,array}
\newcommand{\PROVED}{\textcolor{green!40!black}{\textbf{PROVED}}}
\newcommand{\OPEN}{\textcolor{red!70!black}{\textbf{OPEN}}}
\newcommand{\AUDIT}{\textcolor{orange!70!black}{\textbf{AUDIT}}}
\newtheorem{theorem}{Theorem}
\newtheorem{proposition}{Proposition}
\title{Lambda and determinant control of Witt contractivity (v103)}
\date{13 August 2026}

\begin{document}
\maketitle

\section{Exterior-power criterion}

Let $T$ act on a finite-dimensional semisimple complex vector space $M$.
Give every simple summand a positive metric and the direct sum its
orthogonal metric.

\begin{theorem}[Compound-operator criterion --- \PROVED]
The following are equivalent:
\begin{enumerate}
\item $\|T\|\le1$;
\item $\|\bigwedge^kT\|\le1$ for every $1\le k\le\dim M$;
\item the largest singular value of $T$ is at most one.
\end{enumerate}
If only the top determinant bound
$|\det T|\le1$ is known, none of these conclusions follows when
$\dim M\ge2$.
\end{theorem}

\begin{proof}
The singular values of $\bigwedge^kT$ are all products of $k$ singular
values of $T$; in particular the largest is the product of the $k$ largest
singular values.  Taking $k=1$ proves the equivalence immediately.  For the
last assertion, $T=\operatorname{diag}(R,R^{-1})$ has determinant one and
operator norm $R>1$.
\end{proof}

For a one-dimensional Loewy factor $M_\rho$ no exterior-power ambiguity is
left:
\begin{equation}
 \det(\sqrt n,U_n|M_\rho)=n^{1/2-\rho}.             \tag{1}
\end{equation}
Thus a metric determinant comparison on each simple factor would already
be exactly (8c).

\section{The three determinant constructions in the supplied rows}

The files construct three different objects:
\begin{enumerate}
\item row (a): the top determinant
$\det\mathbf H_t^{\rm cot}(D,E)$ of a finite free code-cotangent module,
with an explicitly selected degree-matching norm;
\item row (b): Haran's lambda characteristic
$\chi_T(\phi_n)=\Phi_n(T)$ and the determinant of the local complex
$[\mathbb Z\xrightarrow{\Phi_n(1)}\mathbb Z]$;
\item row (c): the nuclear trace of the completed Poisson quotient.
\end{enumerate}

They are compared by the scalar identity
\begin{equation}
 -\log\|\det_{\rm tor}\mathbb L_n\|=\Lambda(n),       \tag{2}
\end{equation}
and by the monoidal multiplier label $\Gamma_n\mapsto U_n$.  There is no
displayed morphism which sends an exterior power of the row-(a) cotangent
module or Haran's lambda object to
$\bigwedge^k\operatorname{gr}_LM$ for a row-(c) Loewy quotient.

\begin{proposition}[The existing lambda comparison is not a compound
operator comparison --- \PROVED]
The identity $\chi_T(\phi_n)=\Phi_n(T)$ does not compute the characteristic
polynomial of $\sqrt nU_n$ on any finite Loewy quotient of row (c).
Likewise (2) computes a torsion contact degree, not the determinant in
(1).
\end{proposition}

\begin{proof}
The row-(b) definition applies $\lambda_T$ to the cyclic Witt vector
$\phi_n=V_n\phi_1$ and obtains the cyclotomic polynomial $\Phi_n(T)$.  Its
subsequent derived intersection evaluates that polynomial at $T=1$.  No
row-(c) module $M$, no action matrix on $M$, and no exterior power of that
matrix occurs in the construction.  Formula (2) is the determinant of the
two-term torsion complex resulting from this evaluation.  In contrast,
(1) is the determinant of the scaling action on a spectral quotient.  The
objects have different sources and targets, so the first cannot be
substituted for the second without a new functor.
\end{proof}

\section{The exact lambda-metric promotion}

Let $\mathfrak M_{\rm fin}(E)$ be the finite Loewy subquotients of the odd
object.  A sufficient new structure is a family of positive metric functors
\begin{equation}
 \mathfrak F_M:
 \langle\Gamma_n,\mathbb L_n\rangle_{\lambda,\rm met}
 \longrightarrow\operatorname{End}^{\le1}(\operatorname{gr}_LM)        \tag{3}
\end{equation}
such that
\begin{align}
 \mathfrak F_M(\Gamma_n)&=\sqrt n,U_n|_{\operatorname{gr}_LM},          \tag{4}\\
 \mathfrak F_M(\lambda^k\Gamma_n)&=igwedge^k\mathfrak F_M(\Gamma_n),  \tag{5}\\
 \mathfrak F_M(\Gamma_n^\dagger)&=mathfrak F_M(\Gamma_n)^*.            \tag{6}
\end{align}

\begin{theorem}[Lambda-metric promotion closes row (d) --- \PROVED]
If (3)--(6) are constructed from rows (a)--(c), then (8c), the positive
Hilbert Loewy assembly and row (d) follow.
\end{theorem}

\begin{proof}
The target of (3) consists of contractions, so (4) gives (8c).  Alternatively
(5) and the compound-operator criterion give the same conclusion.  Equation
(6) supplies the Hilbert adjoint.  The v99 trace-preserving Loewy assembly
then identifies the primitive supertrace with the negative trace of a
Hilbert square.
\end{proof}

The actual functor in row (c) is
\[
 \mathcal R:\mathsf{Corr}^{\rm arith}_{\mathbb C}(X)
 \longrightarrow\mathsf{Mult}(\mathcal S_>).
\]
It is proved strict monoidal and complex-linear.  Its target is not a
metric/dagger category, it is not defined on the Loewy factors, and no
claim corresponding to (5) or (6) appears in its contract.  Hence
(3)--(6) remain \OPEN.

The dimension-one test shows that adding all exterior determinants cannot
hide the missing step: on $M_\rho$, (3) itself says
\[
                    |n^{1/2-\rho}|\le1,
\]
which is the desired half-plane inequality.  Therefore the needed
lambda-metric promotion must be constructed by a new positive fiber functor
or boundary frame; it cannot be obtained by reinterpreting (2).

\begin{center}
\begin{tabular}{>{\raggedright\arraybackslash}p{.58\linewidth}
                >{\centering\arraybackslash}p{.15\linewidth}
                >{\raggedright\arraybackslash}p{.17\linewidth}}
\toprule
Statement & Status & Location \\
\midrule
Exterior-power contraction criterion & \PROVED & Theorem 1 \\
Row-(a) and row-(b) determinant constructions & \PROVED & supplied files \\
Identification with compound operators on Loewy factors & false as stated & type audit \\
Metric lambda-functor (3)--(6) & \OPEN & new construction \\
Row (d), condition $G$ & \OPEN & follows from (3)--(6) \\
\bottomrule
\end{tabular}
\end{center}

\end{document}
